mportance ratio. Let
\(R(x,y)\) denote the signed reward used by TOPR and define

\[
\mathcal{T}^{+}
=
\{(x,y):R(x,y)\geq 0\},
\qquad
\mathcal{T}^{-}
=
\{(x,y):R(x,y)<0\}.
\]

We implement canonical TOPR, whose expected update is

\[
\begin{aligned}
g_{\mathrm{TOPR}}(\theta)
={}&
\mathbb{E}_{(x,y)\sim\mu}
\Big[
\mathbf{1}\{(x,y)\in\mathcal{T}^{+}\}
R(x,y)
\nabla_\theta\log\pi_\theta(y\mid x)
\Big]
\\
&+
\mathbb{E}_{(x,y)\sim\mu}
\Big[
\mathbf{1}\{(x,y)\in\mathcal{T}^{-}\}
\min\{\rho_\theta(x,y),1\}
R(x,y)
\nabla_\theta\log\pi_\theta(y\mid x)
\Big].
\end{aligned}
\]

Thus, positive responses receive the SFT-style unit coefficient, while
negative responses are attenuated using a truncated importance ratio.
The behavior-policy log probability required by the ratio is stored
when the frozen response bank is constructed.

\paragraph{Countdown remoteness.}

For a prompt \(x\) and completion
\(y=(y_1,\ldots,y_T)\), learner-relative remoteness is the current mean
completion-token surprisal

\[
D_{\mathrm{CD}}(x,y)
=
-\frac{1}{T}
\sum_{t=1}^{T}
\log
\pi_\theta
\left(
y_t\mid x,y_{<t}
\right).
\]

Only completion tokens are included; prompt tokens and padding positions
are excluded. Mean rather than summed surprisal is used so that response
length does not mechanically determine remoteness. The remoteness score
is detached before applying the shared taper functions.

\paragraph{Frozen response-bank construction.}

The response bank is generated before policy optimization and then held
fixed. Each record contains the puzzle, generated response, verifier
outcome, response length, and the behavior-policy log probability
required by methods that use importance ratios. The same bank and
training split are supplied to every protocol-matched method.

A response is marked successful only when it forms a valid arithmetic
expression, uses the supplied numbers according to the task rules, and
evaluates to the target. Invalid syntax, illegal number use, and
incorrect arithmetic results remain distinct verifier outcomes in the
stored diagnostics, even when they share the same binary training
reward.

\paragraph{Training and checkpoint selection.}

All methods start from the same SFT checkpoint and share the optimizer,
learning-rate schedule, minibatch construction, number of optimizer
steps, validation split, test puzzles, and random seeds. Hyperparameters
are selected using the registered validation rule. The held-out test set
is evaluated only after checkpoint selection.

The registered training and selection configuration is summarized in
Table~\ref{tab:countdown_registered_config}. All entries must be
populated directly from the frozen experiment configuration rather than
reconstructed after observing the test results.

\begin{table}[t]
\centering
\caption{Registered configuration for the Countdown performance experiments.}
\label{tab:countdown_registered_config}

\small
\setlength{\tabcolsep}{8pt}
\renewcommand{\arraystretch}{1.12}

\begin{tabular}{@{}ll@{}}
\toprule
Configuration & Registered value \\
\midrule
Optimizer-step budget
& \textbf{[registered value]} \\
Batch size
& \textbf{[registered value]} \\
Learning-rate schedule
& \textbf{[registered value]} \\
Random seeds
& \textbf{[registered seed set]} \\
AsymRE baseline candidates, $V$
& \textbf{[registered grid]} \\
Near-field threshold candidates, $\tau$
& \textbf{[registered grid]} \\
Remoteness-scale candidates, $c$
& \textbf{[registered grid]} \\
Taper-strength candidates, $\lambda$
& \textbf{[registered grid]} \\
Global-scaling candidates, $\alpha$
& \textbf{[registered grid]} \\
Number of samples for pass@$k$
& \textbf{[registered value]} \\
\bottomrule
\end{tabular}
\end{table}

\paragraph{Verifier-based evaluation.}

Greedy verifier success is the fraction of held-out puzzles solved by
greedy decoding. Pass@\(k\) is the fraction solved by at least one of
the registered \(k\) sampled responses. Valid-expression rate is the
fraction of generated responses that satisfy the parser and number-use
constraints, irrespective of whether they reach the target.

These three metrics are reported separately. In particular, a higher
verifier success rate is not attributed to improved reasoning when it
is accompanied by a deterioration in expression validity.

Task-performance degradation, token- or sequence-support boundary
events, and NaN/Inf numerical failures are audited independently.
Claims about stability or final method ranking use the registered
terminal checkpoint audit rather than finite-step pilot behavior.

\end{document}